\documentclass[letterpaper,12pt]{article}
\usepackage[top=1in,bottom=1in,left=1in,right=1in]{geometry}
\usepackage{fancyhdr}
\usepackage{xcolor}
\usepackage{url}
\usepackage{array}
\usepackage{tabularx}
\usepackage{enumitem}
\usepackage[colorlinks=true,urlcolor=blue,citecolor=blue,linkcolor=blue]{hyperref}

\urlstyle{same}
\setlength{\parindent}{0pt}
\setlength{\parskip}{0pt}
\pagestyle{fancy}
\fancyhf{}
\fancyfoot[C]{\small\thepage}
\renewcommand{\headrulewidth}{0pt}

\setlist[itemize]{topsep=2pt,itemsep=3pt,parsep=0pt,leftmargin=1.4em}

% Hanging-indent paragraph for publication entries
\newcommand{\pub}[1]{%
  \begingroup
  \hangindent=1.5em
  #1\par
  \endgroup
  \vspace{3pt}%
}

\begin{document}

\begin{center}
  {\large\textbf{Marcus W. Beck, Ph.D.}}\\[3pt]
  \textit{Senior Scientist}\\[2pt]
  Tampa Bay Estuary Program\\
  263 13th Ave South, St.\ Petersburg, Florida 33701\\
  Email: \href{mailto:mbeck@tbep.org}{mbeck@tbep.org}\quad Tel: 727-893-2765\\
  GitHub: \href{https://github.com/fawda123}{fawda123}\quad
  ORCID: \href{https://orcid.org/0000-0002-4996-0059}{0000-0002-4996-0059}
\end{center}

\vspace{4pt}
\noindent\textbf{Education}

\vspace{4pt}
\begin{tabularx}{\linewidth}{@{}X p{4.6cm} p{1.4cm} r@{}}
University of Minnesota, Twin Cities, MN     & Conservation Biology & Ph.D. & 2013 \\[1pt]
University of Minnesota, Twin Cities, MN     & Conservation Biology & M.Sc. & 2009 \\[1pt]
University of Florida, Gainesville, FL       & Zoology              & B.Sc. & 2006 \\[1pt]
Santa Fe Community College, Gainesville, FL  & Zoology              & A.A.  & 2004 \\
\end{tabularx}

\vspace{8pt}
\noindent\textbf{Relevant Professional Appointments}

\vspace{2pt}
\textbf{Senior Scientist} \hfill \textbf{Nov. 2019 -- Present}\\
Tampa Bay Estuary Program

\vspace{4pt}
\textbf{Private Contractor} \hfill \textbf{Apr. 2021 -- Present}\\
EnviroDev, LLC

\vspace{2pt}
\textbf{Scientist} \hfill \textbf{Sep. 2017 -- Sep. 2019}\\
Southern California Coastal Water Research Project

\vspace{2pt}
\textbf{Post-Doctorate Researcher} \hfill \textbf{Dec. 2015 -- Sep. 2017}\\
United States Environmental Protection Agency

\vspace{2pt}
\textbf{Post-Doctorate Research Fellow} \hfill \textbf{Jul. 2013 -- Nov. 2015}\\
ORISE, United States Environmental Protection Agency

\vspace{8pt}
\noindent\textbf{Research and Technical Experience Relevant to this Project}

\vspace{4pt}
\begin{itemize}
\item Fifteen years of experience integrating long-term field monitoring, in-situ sensor, and remote-sensing data into operational water quality and seagrass management products for Tampa Bay and other U.S.\ estuaries.
\item Direct collaborator with Dr.\ Frank Muller-Karger (USF College of Marine Science, Scientific PI on this proposal) on cloud-based satellite remote sensing methods for seagrass monitoring (Lizcano-Sandoval et al.\ 2025), developing a working relationship and familiarity with USF CMS seagrass mapping and classification approaches.
\item Completed hands-on training in marine biodiversity data standards (OBIS, Darwin Core) and data handling at the USF/MBON BioData Mobilization Workshop (Jun.\ 2026), taught in part by Dr.\ Ana Carolina Peraltova, Co-Investigator on this proposal.
\item Led synthesis of Tampa Bay water quality and seagrass monitoring data to define seagrass light requirements and depth of colonization (Beck et al.\ 2018) and to identify climate-related water quality conditions suboptimal for seagrass recovery (Beck et al.\ 2024), directly informing the ecological context for validating acoustic seagrass monitoring.
\item Principal developer of tbeptools, an open-source R package supporting dashboard and reporting products used operationally by Tampa Bay Estuary Program and its partners to ingest, QA/QC, analyze, and publicly disseminate estuarine monitoring data (Beck et al.\ 2021).
\item Author of the Tampa Bay Estuary Program's Data Management Standard Operating Procedures (Beck et al.\ 2021), defining intake, quality control, archiving, and public release of estuarine monitoring data, directly transferable to this project's Data Management Plan and IOOS/SECOORA data-sharing requirements.
\item Supported the multi-agency Piney Point wastewater discharge monitoring response (2021), synthesizing water quality, seagrass, and macroalgae data across more than ten partner agencies and delivering reporting products to inform real-time response (Beck et al.\ 2022, 2023).
\item Applies statistical trend models (wqtrends, WRTDStidal) to convert raw monitoring data into actionable indicators of ecosystem condition for restoration managers, directly analogous to the proposed acoustic stressor quantification framework.
\item Extensive experience training technical staff at resource management agencies (NOAA NERRS SWMP, FDEP, TBEP, MassBays, SCCWRP) in the use of new monitoring data tools and dashboards, delivering more than fifteen workshops that build the operational capacity required for working with data.
\end{itemize}

\vspace{4pt}
\noindent\textbf{Selected Publications (most relevant to this project)}

\vspace{4pt}
\pub{Lizcano-Sandoval, L., \textbf{Beck, M.W.}, Scolaro, S., Sherwood, E.T., Muller-Karger, F.E. 2025. Cloud-based satellite remote sensing for enhancing seagrass monitoring and ecosystem management. \textit{ISPRS Journal of Photogrammetry and Remote Sensing}. 227:508--518. {\footnotesize\href{https://doi.org/10.1016/j.isprsjprs.2025.06.034}{10.1016/j.isprsjprs.2025.06.034}}}

\pub{\textbf{Beck, M.W.}, Flaherty-Walia, K.E., Scolaro, S., Burke, M.C., Furman, B.T., Karlen, D.J., Pratt, C., Anastasiou, C.J., Sherwood, E.T. 2024. Hot and fresh: Evidence of climate-related suboptimal water conditions for seagrass in a large Gulf coast estuary. \textit{Estuaries and Coasts}. 47:1475--1497. {\footnotesize\href{https://doi.org/10.1007/s12237-024-01385-0}{10.1007/s12237-024-01385-0}}}

\pub{\textbf{Beck, M.W.}, Hagy, J.D., Le, C. 2018. Quantifying seagrass light requirements using an algorithm to spatially resolve depth of colonization. \textit{Estuaries and Coasts}. 41(2):592--610. {\footnotesize\href{http://dx.doi.org/10.1007/s12237-017-0287-1}{10.1007/s12237-017-0287-1}}}

\pub{\textbf{Beck, M.W.}, Schrandt, M.N., Wessel, M.R., Sherwood, E.T., Raulerson, G.E., Prasad, A.A.B., Best, B.D. 2021. tbeptools: An R package for synthesizing estuarine data for environmental research. \textit{Journal of Open Source Software}. 6(65):3485. {\footnotesize\href{https://doi.org/10.21105/joss.03485}{10.21105/joss.03485}}}

\pub{\textbf{Beck, M.W.}, Wetherill, B., Carr, J. 2023. MassWateR: Improving quality control, analysis, and sharing of water quality data. \textit{PLOS ONE}. 18(11):e0293737. {\footnotesize\href{https://doi.org/10.1371/journal.pone.0293737}{10.1371/journal.pone.0293737}}}

\pub{\textbf{Beck, M.W.}, Burke, M.C., Raulerson, G.E., Scolaro, S., Sherwood, E.T., Whalen, J. 2023. Coordinated monitoring of the Piney Point wastewater discharge into Tampa Bay: Data synthesis and reporting. \textit{Florida Scientist}. 86(2):288--300.}

\pub{\textbf{Beck, M.W.}, de Valpine, P., Murphy, R., Wren, I., Chelsky, A., Foley, M., Senn, D.B. 2022. Multi-scale trend analysis of water quality using error propagation of generalized additive models. \textit{Science of the Total Environment}. 802:149927. {\footnotesize\href{https://doi.org/10.1016/j.scitotenv.2021.149927}{10.1016/j.scitotenv.2021.149927}}}

\vspace{4pt}
\noindent\textbf{Relevant Awards, Certifications, and Professional Memberships}

\vspace{4pt}
\begin{itemize}
\item USEPA Office of Research and Development, Level III Science and Technological Achievement Award, 2022.
\item Coastal and Estuarine Research Federation: Scientific Program Committee Co-Chair, 2025 biennial meeting; Member since 2013.
\item Grant Reviewer: NOAA RESTORE FFO-2025; Florida Sea Grant Graduate Fellowship Competition; Florida FWC Harmful Algal Bloom Grants Program; Florida RESTORE Act Centers of Excellence Program; Maryland Sea Grant; University of Southern California Sea Grant.
\end{itemize}

\vspace{4pt}
\noindent\textbf{Experience Managing Similar Projects}

\vspace{4pt}
\begin{itemize}
\item As Senior Scientist for Tampa Bay Estuary Program, coordinates the synthesis of data to report on the Bay's long-term water quality and seagrass monitoring across partner agencies, translating monitoring data into the annual Bay Health report card and management actions reviewed by TBEP's Boards and Technical Advisory Committee.
\item Supported the multi-agency Piney Point monitoring response (2021): coordinated data aggregation, QA/QC, and synthesis across more than ten partner agencies and delivered public data products to support real-time management decisions (Beck et al.\ 2022 \textit{Marine Pollution Bulletin}; Beck et al.\ 2023 \textit{Florida Scientist}).
\item Co-Principal investigator and developer for MassWateR, a data management and QA/QC software product developed under contract with the Massachusetts Bays National Estuary Program; supported in-person and remote trainings to volunteer and professional monitoring programs across Massachusetts.
\item Co-Chair, Scientific Program Committee for the 2025 Coastal and Estuarine Research Federation biennial meeting, coordinating conference science themes and abstract review across a multi-institution program committee.
\item Managed independent software and data-analysis contracts (as EnviroDev, LLC) for multiple federal, state, and non-profit partners, including the San Francisco Estuary Institute, Penn State University, the Southern California Coastal Water Research Project, and the Choctawhatchee Basin Alliance, each with defined deliverables, budgets, and timelines.
\end{itemize}

\end{document}
